%======================================================================
%  Companion to Chapter 5 -- Plasticity: integration algorithms
%======================================================================
\chapter{Plasticity: integration algorithms}\label{cc:ch5}

\framepart{Outline}

Chapter~\ref{B-ch:pl-num} of the notes builds an elastoplastic computation
from two pieces: a return map at each Gauss point (Box~5.2) and an
equilibrium loop on the nodal unknowns (Box~5.3). The Python module
\texttt{ci\_core} writes the return map point by point. Gibiane works on whole
fields, and the same algorithm becomes a short sequence of operations on the
Gauss-point stress field: no loop over the points, the scalar Newton
iteration of Box~5.2 done everywhere at once. This chapter writes that return
map, drives it at a material point, then puts it inside an equilibrium loop
for the thick-walled cylinder and compares it with \op{pasapas}.

You should then be able to read a constitutive integration algorithm as an
operation on fields, program it in gibiane, and see the cost of the iteration
matrix: the elastic stiffness (linear convergence) against the consistent
tangent of the notes (quadratic convergence).

\section{Where Cast3M enters Chapter 5}\label{cc:sec5-map}
\index{Cast3M!Chapter 5}%

\begin{gbox}{Chapter 5 of the notes and Cast3M}\label{cc:tab5}
\small
\begin{tabular}{@{}p{42mm}@{\hspace{3mm}}p{98mm}@{}}
Radial return, Box~5.2, Sec.~\ref{B-sec:pn-j2};
Ex.~\ref{B-exo:pn-radial}, \ref{B-exo:pn-shear}
 & \textbf{Exercise \ref{exo:cc-rrj2}}: the procedure \texttt{rrj2} on Gauss-point fields
   (linear or Voce isotropic, linear kinematic hardening), driven on one
   element with a prescribed homogeneous strain.\\[1mm]
Equilibrium, Box~5.3, Sec.~\ref{B-sec:pn-newton}; Ex.~\ref{B-exo:pn-fem},
\ref{B-exo:pn-tangent}
 & \textbf{Exercise \ref{exo:cc-cylinder}}: \texttt{rrj2} inside an equilibrium loop with the
   elastic stiffness; the cylinder at $0.8p_L$, $0.95p_L$ and with Voce
   hardening; against \op{pasapas}.\\[1mm]
Viscoplasticity, Sec.~\ref{B-sec:pn-vp}
 & In \texttt{rrj2}: the entries \texttt{ETA} and \texttt{DT} of the material
   table give the Perzyna update~\eqref{B-eq:pn-perzyna-j2}. To do: the
   relaxation of Exercise~\ref{B-exo:pl-perzyna}, Duvaut--Lions.\\[1mm]
Corners, Box~5.4; Ex.~\ref{B-exo:pn-tresca}
 & To do: Tresca on the principal stress fields (\op{prin}), the active set
   decided by masks, corner test $2f_b>f_a$.\\[1mm]
Plane stress, Box~5.5; Ex.~\ref{B-exo:pn-ps}
 & To do: the projected return on fields; the scalar equation for $\dg$ solved
   by Newton at all points, as in \texttt{rrj2}.\\[1mm]
Consistent tangent, Sec.~\ref{B-ssec:pn-calg}
 & To do: an iteration matrix built from a Gauss-point modulus field
   (\op{mate} with fields), or \op{pasapas} with its tangent option; compare
   the iteration counts with Exercise~\ref{exo:cc-cylinder}.\\
\end{tabular}
\end{gbox}

\section{The radial return on fields}\label{cc:sec5-rrj2}
\index{radial return!on fields}%
In Box~5.2 the multiplier is $\dg$ and the norm is $\norm{\teta}$. Gibiane
has the von Mises norm \op{vmis}, $q=\sqrt{3/2}\,\norm{\dev\teta}$, so the
procedure uses $q$ and the increment of the cumulated plastic strain
$\Delta\alpha=\sqrt{2/3}\,\dg$, called \texttt{dp}. In these variables
Box~5.2 reads as follows.

\begin{gbox}{Box C5.1\quad Radial return on Gauss-point fields}\label{cc:box-rrj2}
\index{radial return!Box C5.1}%
\begin{enumerate}[label=\textbf{\arabic*.},leftmargin=1.6em]
  \item \textbf{Predictor.} With the displacement $\mat{u}_n$ of the
  converged state, $\sig^\trial=\sig_n+\bbC:\eps(\mat{u}-\mat{u}_n)$, one
  call of \op{sigm}; $\teta^\trial=\sig^\trial-\tbeta_n$ and
  $q^\trial=\sqrt{3/2}\,\norm{\dev\teta^\trial}$, one call of \op{vmis}.
  \item \textbf{Mask.} $f^\trial=q^\trial-R(\alpha_n)$; the field
  \texttt{pla} is 1 where $f^\trial>0$ and 0 elsewhere (\op{masq}).
  \item \textbf{Multiplier.} From $\Delta\alpha=0$, at all points at once,
  \[
    g=\texttt{pla}\,\bigl(q^\trial-a\,\Delta\alpha
      -R(\alpha_n+\Delta\alpha)\bigr),\qquad
    \Delta\alpha\leftarrow\Delta\alpha+\frac{g}{a+R'(\alpha_n+\Delta\alpha)},
  \]
  with $a=3\mu+H+\tfrac32\eta/\Delta t$, until $\max\abs{g}<10^{-10}\sY$: one iteration for linear hardening,
  a few for Voce. $\eta=0$ is the rate-independent case.
  \item \textbf{Update.} With $\dev\teta^\trial$ built component by
  component,
  \[
    \sig_\np=\sig^\trial-3\mu\,\Delta\alpha\,\frac{\dev\teta^\trial}{q^\trial},\quad
    \tbeta_\np=\tbeta_n+H\,\Delta\alpha\,\frac{\dev\teta^\trial}{q^\trial},\quad
    \alpha_\np=\alpha_n+\Delta\alpha .
  \]
\end{enumerate}
\end{gbox}

\paragraph{Why it is Box~5.2.}
With $\Delta\alpha=\sqrt{2/3}\,\dg$, $q=\sqrt{3/2}\norm{\teta}$ and
$\tn=\sqrt{3/2}\,\dev\teta/q$: $2\mu\dg\,\tn=3\mu\,\Delta\alpha\,\dev\teta/q$,
$\tfrac23H\dg\,\tn=H\,\Delta\alpha\,\dev\teta/q$, and $g(\dg)=0$ of Box~5.2
multiplied by $\sqrt{3/2}$ is the equation of step~3. For linear hardening
$\Delta\alpha=f^\trial/(3\mu+K+H)$, which is $\dg=f/(2\mu+\tfrac23(K+H))$ of
the notes. The Perzyna term $\tfrac32\eta/\Delta t$ is
equation~\eqref{B-eq:pn-perzyna-j2} in the same variables.

\begin{castemcom}
\textbf{Predictor and mask.} \texttt{sgt} is the trial stress, computed by
the caller from the converged state. \op{vmis} gives the norm $q^\trial$ as a
field with one component, \texttt{SCAL}.

\op{masq} with the threshold $-1$ gives a field of ones on the same Gauss
points, used to turn numbers into fields; with the threshold 0 on
$f^\trial$ it gives the plastic points.
\end{castemcom}
\begin{castemlines}
\begin{verbatim}
eta = sigtr - bet;
qtr = vmis mo eta;
un  = masq qtr 'SUPERIEUR' -1.;
rr0 rp0 = hard al un tm;
ftr = qtr - rr0;
pla = masq ftr 'SUPERIEUR' 0.;
\end{verbatim}
\end{castemlines}
\begin{castemcom}
\textbf{The scalar Newton iteration, everywhere at once.} \texttt{hard}
returns the fields $R$ and $R'$ (linear or Voce). The residual \texttt{gg}
is zero at the elastic points, where \texttt{dp} stays zero. The loop stops
on the largest residual over the mesh.
\end{castemcom}
\begin{castemlines}
\begin{verbatim}
dp = 0. * un;
repeter bnw 30;
  rr rp = hard (al + dp) un tm;
  gg = pla * (qtr - (a3 * dp) - rr);
  si ((maxi (abs gg)) < (1.e-10 * sy));
    quitter bnw;
  finsi;
  dp = dp + (gg / ((a3 * un) + rp));
fin bnw;
\end{verbatim}
\end{castemlines}
\begin{castemcom}
\textbf{Update along the deviator.} \texttt{devia} builds
$\dev\teta^\trial$ component by component, \texttt{smul} multiplies each
component by a scalar field. Both use \op{exco} to rename a component to
\texttt{SCAL} and back, and \op{et} to join the components into one field.
$q$ is shifted by $10^{-12}\sY$ to avoid $0/0$ at the points where
$\Delta\alpha=0$.
\end{castemcom}
\begin{castemlines}
\begin{verbatim}
qs  = qtr + ((1.e-12 * sy) * un);
ed  = devia eta ldia loff;
sig = sigtr - (smul
      ((3. * mu) * (dp / qs)) ed lcp);
bet2 = bet + (smul
      (hh * (dp / qs)) ed lcp);
al2  = al + dp;
\end{verbatim}
\end{castemlines}

\section{The equilibrium loop}\label{cc:sec5-loop}
\index{initial-strain iteration!Cast3M}%
\begin{castemcom}
Box~5.3 with the elastic stiffness \texttt{kk} as iteration matrix: it is
assembled once, outside the loops. This is the initial-strain iteration of
Box~6.1 of the notes (Section~\ref{B-sec:cy-orders}), a modified Newton
method (Section~\ref{B-sec:it-newton}): each iteration is one solve, the
convergence is linear.

The local stage always starts from the converged state
(\texttt{un, sn, bn, an}), never from the previous iterate. The residual and
its test are those of the hand-written loop of the Cast3M introduction:
\texttt{aa * uu} adds the multiplier forces and the constraint residual.
\end{castemcom}
\begin{castemlines}
\begin{verbatim}
uu = un + (reso (kk et aa) (ff - fo));
repeter bequ 2000;
  nit = nit + 1;
  sgt = sn + (sigm mo ma (uu - un));
  s b a dp ft = rrj2 mo sgt bn an tm;
  rr = ff - (bsig mo s) - (aa * uu);
  res = maxi (abs (exco rr lfor));
  si (res < (tol * (maxi
            (abs (exco ff lfor)))));
    quitter bequ;
  finsi;
  uu = uu + (reso (kk et aa) rr);
fin bequ;
\end{verbatim}
\end{castemlines}

\framepart{Summary}

\begin{gbox*}
\begin{itemize}
\item \textbf{A return map is an operation on fields}: \op{sigm} gives the
  predictor, \op{vmis} the norm, \op{masq} the active set, a scalar Newton
  iteration runs at all Gauss points at once, and the update is a product of
  a scalar field by the deviator.
\item \textbf{The variables of Box~5.2 change, not the algorithm}:
  $\Delta\alpha=\sqrt{2/3}\,\dg$ and $q=\sqrt{3/2}\norm{\teta}$.
\item \textbf{Start every iteration from the converged state}: the
  procedure receives $\sig_n$, $\tbeta_n$, $\alpha_n$ and the total increment
  since $t_n$.
\item \textbf{The iteration matrix decides the cost}: the elastic stiffness
  gives linear convergence and many cheap iterations; Newton with the
  consistent tangent gives a few expensive ones.
\end{itemize}
\end{gbox*}

\framepart{Exercises}

\exercise{The radial return at a material point}\label{exo:cc-rrj2}
\index{radial return!Cast3M}%
Program Box~C5.1 as a procedure on Gauss-point fields. Drive it on one
\texttt{QUA4} element in plane strain whose nodes all follow $\vu=\eps\cdot\vx$:
the strain is homogeneous and no equilibrium iteration is needed. With the data
of Exercise~\ref{B-exo:pn-shear} ($E=200$\,GPa, $\nu=0.3$, $\sY=250$\,MPa,
$K=2$\,GPa): (a)~simple shear $\gamma=0.01$ in one step and in 1000 steps;
(b)~shear to $\varepsilon_{12}=0.002$, then $\varepsilon_{11}$ from 0 to 0.004 at
fixed shear in $n=1$, 10, 100, 1000 steps. Plot the error on $\sigma_{12}$
against $n$.
\begin{solution}
(a)~$\sigma_{12}=149.706775$\,MPa in one step and in 1000: the path is
proportional and the return map is exact. (b)~The reference (200\,000 steps)
is $\sigma_{11}=834.8243$, $\sigma_{12}=26.0140$\,MPa,
$\alpha=3.114637\times10^{-3}$; the error on $\sigma_{12}$ is 30.07, 4.941,
0.5281, 0.05293\,MPa: first order, the normal turns during the step
(\texttt{exo\_shear.py}). The gibiane numbers must agree to the printed digits:
the procedure is the same algorithm on a different data structure.
\untested{\texttt{cast3m/ch5/return\_map\_point.dgibi}}
\lstinputlisting[style=gibstyle,firstline=31]{cast3m/ch5/return_map_point.dgibi}
\end{solution}

\exercise{The thick-walled cylinder, return map and equilibrium in gibiane}\label{exo:cc-cylinder}
\index{thick-walled cylinder!Cast3M}%
Put the procedure of Exercise~\ref{exo:cc-rrj2} in the equilibrium loop of
Section~\ref{cc:sec5-loop} for the cylinder of Exercise~\ref{B-exo:pn-fem}
($a=100$, $b=200$\,mm, plane strain, axisymmetric slice of 200
\texttt{QUA4}). (a)~Perfect plasticity, $\sY=250$\,MPa: 20 steps to $0.8p_L$
and to $0.95p_L$; the bore displacement, the radius of the plastic zone, the
number of iterations. (b)~Voce hardening ($\sigma_\infty=600$\,MPa, $\delta=30$),
$p=320$\,MPa in 10 steps. (c)~The same computation (a) with \op{pasapas}.
\begin{solution}
The Python code of the notes (Newton with the consistent tangent, 200 linear
elements) gives $p_L=200.09$\,MPa; at $0.8p_L$, $u(a)=0.18305$\,mm and
$c=131.0$\,mm (closed form 130.8) in 60 iterations; at $0.95p_L$,
$u(a)=0.30536$\,mm and $c=165.5$\,mm (164.0) in 70 iterations. With Voce
hardening and 80 elements, $u(a)=3.4526$\,mm in 47 Newton iterations
(Exercise~\ref{B-exo:pn-tangent}). The gibiane loop should give the same
displacements to the mesh difference and the plastic radius to one Gauss-point
spacing (0.5\,mm), with many more iterations: the elastic stiffness converges
linearly, with a factor close to 1 when most of the wall is plastic, so that
the count grows quickly near $p_L$. This is the price of a matrix built once,
and the reason for the consistent tangent. \op{pasapas} gives the reference
of~(c).
The listing starts after the procedures \texttt{smul}, \texttt{devia},
\texttt{hard} and \texttt{rrj2}, printed with Exercise~\ref{exo:cc-rrj2}.
\untested{\texttt{cast3m/ch5/cylinder\_radial\_return.dgibi}}
\lstinputlisting[style=gibstyle,firstline=155]{cast3m/ch5/cylinder_radial_return.dgibi}
\end{solution}

\begin{chapterbib}{9}
\bibitem{cc-simo} J.C.~Simo and T.J.R.~Hughes, \emph{Computational
  Inelasticity}, Springer, 1998.
\end{chapterbib}
